\cotefile{156-9}
\scanpage{112}
\mspage{NA 20}
% archive sheet number: 111
% Typescript upside down in the scan; transcribed in its reading direction.
% One diagonal cancellation stroke crosses the entire body of the typescript.
% Nested deletion marks retain the separate local cancellations.
\deleted{(que je qualifie de ``foncteur de Mebkhout'' dans le texte) figurant dans la conjecture en question, m'a assuré son entière bonne foi, vu qu'il n'aurait pas eu connaissance de cet appendice à l'exposé de Ramis, avant que je ne lui communique la lettre de Schapira en faisant état, au mois de janvier dernier.}

\deleted{Page 951}

\deleted{($^{**}$) (6 mai 1986) C'est là du moins la ``version Mebkhout'' des choses, contredite par d'autres témoignages qui m'ont paru convainquants. Je n'ai pas de doute à présent que l'idée des foncteurs $m$ et $\unsure{M|\mathbb{M}}$ (formant la fameuse ``correspondance'' dite de Riemann-Hilbert), dans le cadre des catégories dérivées, est due à Kashiwara. \deleted{Mais \unsure{personnellement}, \unsure{jusqu'en décembre 1979},} cette conjecture \deleted{\unsure{a été regardée}} semble avoir été regardée, par ceux qui en avaient entendu parler tout comme par Kashiwara et Mebkhout, comme une curiosité mathématique sans grande conséquence, du moins jusqu' en septembre 1979.}

\deleted{Page 952}

\deleted{($^{*}$) (6 mai 1986) C'est là du moins la version Mebkhout, que j'avais prise pour argent comptant. J'en doute fortement à présent, et surtout (comme je le souligne dans de précédentes notes de bas de page, datées de ce jour) qu'il ait eu avant fin 1979 ou courant 1980 l'intuition du rôle crucial que cette équivalence de catégories (conjecturée d'abord par Kashiwara, selon toute apparence) allait avoir à jouer dans une nouvelle théorie de coefficients. Il convient donc de lire également avec un grain \deleted{\unsure{à}} de sel la note de b. de p. qui suit (du 5 mai l'an dernier).}

\deleted{Complément à la note de bas de page ($^{**}$)}

\deleted{($^{\unsure{x|\times}}$) (6 mai 1986) Pour une rectification au sujet de cette version, inexacte, voir les précisions datées de ce même jour, dans les notes de b. de p. ($^{*}$) p. 949 et ($^{*}$) p. 950: \deleted{\unsure{l'an dernier}} la conjecture pertinente est mentionnée dès janvier 1978 dans un texte de la plume de Ramis, qui l'attribue à Kashiwara.}

\deleted{Page 953}

\deleted{($^{*}$) (6 mai 1986) Pour des réserves vis-à-vis de cette terminologie ``foncteur de Mebkhout'', voir note de b. de p. ($^{\unsure{***|**}}$) p. 949. On pourrait peut-être aussi l'appeler ``foncteur de Kashiwara-Mebkhout'', vu le rôle important \deleted{\unsure{pris}} joué par Mebkhout, aussi bien pour faire réaliser son importance (via sa philosophie de dualité, notamment), que pour \unsure{le|la} \deleted{\unsure{l'immédiat}}}

% A small black mark below the last line has no recoverable letter shape.
